% THREE-NODE-MIXED-20261016: audited appendix integration
\section*{Three-node mixed-deadline CPU experiment}
\label{sec:three-node-mixed}
This is a new three-node experiment for heterogeneous request deadlines, distinct from the earlier single-host measurements retained in Table~S4. It complements the resource-diversity experiment in main-text Fig.~5. NJ-k1 is the cloud node; NJ-k2/NJ-k3 are request origins with local execution capacity. The three Debian~13 KVM virtual machines expose 16 logical CPUs each; NJ-k1/NJ-k2 use Xeon Platinum 8163 processors and NJ-k3 uses a faster Xeon Gold 6442Y. The unshaped LAN and fixed warm deployment do not represent WAN, mobility, accelerator or autoscaling conditions.

\paragraph{Shared workers and resource scope.}
Every system has four prestarted serial Go execution slots allocated 2/1/1. Each application worker has a two-CPU quota, a 512-MiB memory limit and Go parallelism one. Bare, Ray and KubeEdge use cloud CPU masks 0,1 and 2,3; native Tide/PAIX uses 0,2 and 1,3. Both use four cloud logical CPUs, but the grouping is not identical. End-side workers use CPUs 0,1; new control processes and clients use CPUs 8--15. Framework control/runtime processes are outside the application-worker limits, and their total budgets are not identical. All four deployments remain resident, with traffic applied to one system at a time. All requests use the same Go kernel: 65,536 bytes and 100 SHA-256 iterations, with independently checked returned checksums and processed-byte counts.

\paragraph{Fresh calibration and frozen settings.}
Twenty sequential calibration samples after three excluded warmups per node give Bare service-time medians 20.699456/20.7434155/4.1519125~ms and complete local native Tide path medians 33.2802745/31.8464075/9.0587625~ms for NJ-k1/NJ-k2/NJ-k3. The latter are rounded up to execution estimates 34/32/10~ms. With $s_i$ denoting Bare service time in seconds, $C_{\rm CPU}=2/s_1+1/s_2+1/s_3=385.6818254$ requests/s. The nominal 30\%/80\%/120\% targets round to 116/309/463 total requests/s across both origins. Capacity is a pre-outcome service-time reference, not a prediction of framework throughput. Slot-count round robin does not compensate for NJ-k3's faster processor.

Tide retains equal score weights $(1,1,1)$, a 100-ms latency scale, 5000-ms maximum state age, 1-ms configured per-node network estimates and a 24-hour Stub TTL. The native request function names \texttt{cpu-qos100} and \texttt{cpu-qos500} map to actual 100-ms and 500-ms admission deadlines through \texttt{qos\_by\_function\_ms}. Native admission stamps \texttt{SubmitTime} when the receiving ingress accepts the request; client QoS scoring instead uses actual origin submission-to-return latency, including ingress and return transport. Thus the configured native deadline and the measured client deadline have distinct starting boundaries. Both classes share the same four workers; there are no class-specific worker pools. In the baseline deployments, the Go worker does not use the deadline field for admission; the same classes set client-side evaluation deadlines, with the existing framework queues and routing retained. Ray Serve~2.58.0 retains four fixed replicas, one ongoing request each and native routing with local-node preference. KubeEdge~1.23.1 uses real EdgeCore workers on both end-side VMs and Kubernetes-managed cloud workers. Bare/KubeEdge use four-slot round robin with an origin phase offset. Ray's local HTTP forwarding and KubeEdge's explicit ingress adapters are included in end-to-end measurements. This comparison measures integrated admission, queueing, routing and framework effects and does not isolate an individual scheduling mechanism.

\paragraph{Class-balanced traces and measurement.}
Only deterministic uniform arrivals are tested. The frozen trace divides offers equally between origins and globally equally between 100-ms and 500-ms classes. Each origin's class counts differ by at most one. Classes are balanced over eight requests per origin so each round-robin slot receives both classes; they are not a strictly alternating global deadline sequence. The same class-bearing workload bodies and trace are used for all systems. Each offer is attempted once, without retries, a client concurrency cap or a serial client queue.

Five paired repetitions randomize the three load cells and four systems within each cell. The complete campaign has 60 treatments, 120 origin runs and 177,600 offers. Each treatment has a 10-s submission/observation window. Latency is measured from actual origin-local monotonic submission to return, with an inclusive class deadline. On-time goodput counts checksum-valid returns that meet that deadline and return inside the fixed half-open window, divided by 10~s. Cohort yield divides all timely successful returns (including drain) by all class offers; failures and timeouts remain in denominators. Planned-arrival-to-return timing, replay lateness and successful latency including drain remain in the audit. Successful-latency and returned-destination statistics exclude failed offers and cannot describe their destinations. Wall-clock mapping supports ordering, not cross-host latency or queue-time inference.

The independent raw audit verifies every frozen offer, class/function body, origin commit and checksum, and all 24 aggregate class cells. It retains 51,308 failed offers. Treatments wait for responses and observed quiescence before proceeding; bounded CPU/socket/thread checks do not prove an exact empty native admission queue. Intervals are pointwise, unclipped Student-$t$ 95\% intervals over five treatment-level class results, each combining both origins (four degrees of freedom); they do not establish simultaneous coverage. One three-VM allocation, only uniform arrivals and five repetitions limit generalization. Table~\ref{tab:three-node-mixed} reports every system, class and load, including unfavorable outcomes. Earlier single-host results remain separately reported in Table~S4; differences between experiments cannot be attributed to topology alone.

\paragraph{Observed tradeoff.}
At the 80\% target (309 total requests/s), 100/500-ms class goodput is Tide 3.06/147.10, Bare 39.76/47.16, Ray Serve 2.46/8.58, KubeEdge 39.78/47.10 requests/s. At the 120\% target (463 total requests/s), 100/500-ms class goodput is Tide 1.70/147.50, Bare 58.64/63.34, Ray Serve 1.24/5.86, KubeEdge 58.66/63.42 requests/s. The higher-load advantage is confined to the 500-ms class; Tide underperforms Bare/KubeEdge for the 100-ms class. These results do not establish a general advantage across deadline classes.

\begin{figure}[htbp]
\centering
\includegraphics[width=\linewidth]{revision-figs/evaluation/figS7_cluster_mixed.pdf}
\caption{Three-node mixed-deadline CPU measurements at nominal 30\%/80\%/120\% of freshly calibrated CPU capacity (116/309/463 total requests/s). The 100-ms and 500-ms classes share four warm serial workers per system, allocated 2/1/1. Bars show means and error bars show pointwise run-level 95\% confidence intervals from five paired 10-s runs. Dotted references show each class's offered rate, half the total rate; failures remain in all-offer cohort denominators. The comparison includes integrated admission, queueing and framework paths.}
\label{fig:three-node-mixed}
\end{figure}

\begin{table}[htbp]
\centering\small
\caption{Complete three-node mixed-deadline CPU on-time goodput (requests/s). Each cell is the five-run mean $\pm$ its run-level 95\% confidence half-width; all 24 class/load/system cells are retained. Percentages denote nominal calibrated-capacity targets before rounding.}
\label{tab:three-node-mixed}
\setlength{\tabcolsep}{4pt}
\begin{tabular*}{\linewidth}{@{\extracolsep{\fill}}llrrrr@{}}
\toprule
\shortstack{Class\\(ms)} & \shortstack{Load (\% of $C_{\rm CPU}$)\\(total requests/s)} & Tide & Bare & Ray Serve & KubeEdge\\
\midrule
100 & 30 [116] & $57.84 \pm 0.07$ & $57.90 \pm 0.00$ & $55.52 \pm 6.61$ & $57.90 \pm 0.00$\\
100 & 80 [309] & $3.06 \pm 0.14$ & $39.76 \pm 0.07$ & $2.46 \pm 0.88$ & $39.78 \pm 0.06$\\
100 & 120 [463] & $1.70 \pm 0.12$ & $58.64 \pm 0.26$ & $1.24 \pm 0.27$ & $58.66 \pm 0.11$\\
\midrule
500 & 30 [116] & $57.76 \pm 0.07$ & $58.00 \pm 0.00$ & $55.82 \pm 5.78$ & $58.00 \pm 0.00$\\
500 & 80 [309] & $147.10 \pm 1.52$ & $47.16 \pm 0.19$ & $8.58 \pm 2.02$ & $47.10 \pm 0.36$\\
500 & 120 [463] & $147.50 \pm 2.79$ & $63.34 \pm 0.30$ & $5.86 \pm 0.53$ & $63.42 \pm 0.16$\\
\bottomrule
\end{tabular*}
\end{table}
